\section{Final contour moves, collisions, and completion of the proof}
\label{sec:final}

The two $u$-moves in this section are performed only after
\Cref{prop:full-kernel-continuation}.  They are distinct: first we move the
full diagonal boundary, and then the complete three-permutation expression.

\subsection{The diagonal-boundary \texorpdfstring{$u$}{u}-move}

The ordinary $u=0$ residue of \eqref{eq:diagonal-boundary-Y} is
\begin{align}
 \mathcal R_{\mathrm{diag},0}(h,k)
 ={}&\int_{\R}w(t)\bigg\{
 \frac{1}{2\pi i}\int_{(\sigma_s)}\frac{G(s)}s
 g_{\beta,\gamma}(s,t)\notag\\
 &\hspace{31mm}\times
 Z_{\alpha,\beta+s,\gamma+s}(h,k)\dd s\notag\\
 &\quad+X_{\beta,\gamma}(t)
 \frac{1}{2\pi i}\int_{(\sigma_s)}\frac{G(s)}s
 g_{-\gamma,-\beta}(s,t)\notag\\
 &\hspace{31mm}\times
 Z_{\alpha,-\gamma+s,-\beta+s}(h,k)\dd s
 \bigg\}\dd t.
 \label{eq:diagonal-boundary-zero}
\end{align}
Move each $u$-line in \eqref{eq:diagonal-boundary-Y} to
\begin{equation}\label{eq:boundary-u-line}
                  \Re u=-\sigma_u,\qquad \sigma_u=M+\frac12.
\end{equation}
On $\Re s=\sigma_s$, the only $u$-dependent arithmetic poles are
\begin{equation}\label{eq:boundary-arithmetic-poles}
                  u=-\alpha-\gamma-s,\qquad
                  u=\beta-\alpha-s.
\end{equation}
Both have real part $1/2-\eps_s+O((\log T)^{-1})>c_u$, so they lie to the
right of the original line.  The pole at $u=0$ gives
\eqref{eq:diagonal-boundary-zero}; the potential poles
$u=-1,\ldots,-M$ vanish by \eqref{eq:gamma-RM}.

For the new-line estimate, the finite Euler factor in
\eqref{eq:Z-euler} may be written
\begin{equation}\label{eq:Bnu-form}
 B_\nu(q,v)=X^\nu+(1-X)\sum_{j=0}^{\nu-1}X^jY_0^{\nu-j},
 \qquad X=p^{-q},\quad Y_0=p^{-v}.
\end{equation}
This follows by summing the geometric tail in
\[
 (1-p^{-q})(1-p^{-v})\sum_{i+j\geq\nu}p^{-qi-vj}.
\]
On \eqref{eq:boundary-u-line}, \eqref{eq:Bnu-form} and vertical-strip zeta
bounds give
\begin{align}
 Z_{\alpha+u,\beta+s,\gamma+s}(h,k),\quad
 Z_{\alpha+u,-\gamma+s,-\beta+s}(h,k)
 \ll_M{}&(\log T)^{C_M}(hk)^\eps
 h^{\sigma_u+O(\eps_s)+\eps}
 \notag\\
 &\times(1+|u|+|s|)^{C_M}.
 \label{eq:boundary-Z-bound}
\end{align}
Together with $g(s,t)\ll T^{\sigma_s}(1+|s|)^{C_M}$, the gamma decay in
$u$, and the Gaussian decay in $s$, the new lines after the outer coefficient
sum are
\begin{equation}\label{eq:boundary-new-line}
 \ll T^{1+\sigma_s+\eps}Y^{-\sigma_u}
 H^{\sigma_u+O(\eps_s)+\eps}(HK)^{1/2}.
\end{equation}
Since $H\leq T^{2(1-\eta)/3}$ and $Y=T^{B_0}$ with $B_0>2$, $M$ can be
chosen so that \eqref{eq:boundary-new-line} is $O_D(T^{-D})$.  At $u=0$ the
Euler arguments on $\Re s=\sigma_s$ remain a fixed distance from their poles,
and hence
\begin{equation}\label{eq:boundary-zero-bound}
 \sum_{h\leq H}\sum_{k\leq K}
 \frac{|a_hb_k|}{\sqrt{hk}}
 |\mathcal R_{\mathrm{diag},0}(h,k)|
 \ll_\eps T^{1/2+\eps}(HK)^{1/2}.
\end{equation}
Consequently
\begin{equation}\label{eq:boundary-move-result}
 \mathscr D_T
 =\sum_{h\leq H}\sum_{k\leq K}
 \frac{a_h\overline{b_k}}{\sqrt{hk}}
 \mathcal R_{\mathrm{diag},0}(h,k)+O_D(T^{-D}).
\end{equation}

\subsection{The permutation \texorpdfstring{$u$}{u}-move}

\begin{lemma}[Left-line bound]\label{lem:P-left-line}
For fixed $\sigma>0$ and
$\Re\lambda=-\sigma+O((\log T)^{-1})$,
\begin{equation}\label{eq:P-left-bound}
 Q(u)\Pcal_\lambda(h,k;t)
 \ll_{\sigma,M}(\log T)^{C_M}(hk)^\eps
 (T^\sigma+h^\sigma)(1+|\Im u|)^{C_M}.
\end{equation}
After multiplication by $\Gamma(u)$, this is integrable on the complete
vertical line.
\end{lemma}

\begin{proof}
Retain the complete sum $\Pcal_\lambda$ before estimating it.  In the
collision coordinates
\[
 x=\lambda+\gamma,\qquad y=-\beta-\gamma,
 \qquad z=\lambda-\beta=x+y,
\]
the line $\Re\lambda=-\sigma+O((\log T)^{-1})$ keeps $x$ and $z$ a fixed
positive distance from zero.  Thus only $y=0$ can meet this line.

Away from $|y|<1/(2\log T)$, use \eqref{eq:Bnu-form} term by term.  If the
first shift of $Z$ is $\lambda$ and its third shift is $\gamma$ or
$-\beta$, the local product is at worst $h_0^{\sigma+\eps}$.  If the third
shift is $-\lambda$, both Euler arguments have real part
$1+\sigma+O((\log T)^{-1})$, and the local factor cancels the apparent
exterior $h_0^{1/2+\sigma}$, leaving $h_0^{-1/2+\eps}$.  Stirling gives
$X_{\lambda,\gamma}(t)\ll T^{\sigma+\eps}$; the other exact $X$ factor has
size $T^{O((\log T)^{-1})}$.  Standard vertical-strip zeta bounds then give
\eqref{eq:P-left-bound}, with the possible distance to $y=0$ absorbed by a
fixed power of $\log T$.

In the remaining tube $|y|<1/(2\log T)$, apply the joint removability of the
\emph{complete} permutation sum from the logically prior
\Cref{prop:P-holomorphic} before taking absolute values.  With $\gamma$
fixed, write
$\beta=-\gamma-w/\log T$.  The maximum principle on $|w|\leq1/2$ bounds
every value in the tube by the already established bound on $|w|=1/2$.
That circle changes only the logarithmic factor, and $x,z$ remain uniformly
separated from zero.  This proves \eqref{eq:P-left-bound} throughout the
tube, including the collision, without estimating either meromorphic
summand there.  The vertical-strip zeta bounds and the uniform form of
Stirling's formula used above hold for every $\Im u$; the complementary
compact ranges of the gamma arguments are bounded directly.  Thus
\eqref{eq:P-left-bound} holds on the complete vertical line, where
$\Gamma(u)$ dominates the displayed polynomial in $|\Im u|$.
\end{proof}

Set $\sigma=M+1/2$ and move the common $u$-line in
\eqref{eq:R-T} from $c_u$ to $-\sigma$.  The complete residue list is as
follows.
\begin{enumerate}[label=\textup{(\roman*)},leftmargin=8mm]
\item At $u=0$, $\Gamma(u)$ has residue one and $Q(0)=1$; this gives
$\Pcal_\alpha$.
\item The gamma residues at $u=-j$, $1\leq j\leq M$, vanish because
$R_M(-j)=0$.
\item The arithmetic collision locations $u=-\alpha-\gamma$ and
$u=\beta-\alpha$ are removable in the complete $\Pcal$ by
the logically prior \Cref{prop:P-holomorphic}.
\item The poles $u=1/2+2n-\alpha-it$ of
$X_{\alpha+u,\gamma}(t)$ lie to the right of $c_u$ and are not crossed.
The $n=0$ pole encountered in the preliminary move from $\Re u=1$ to $c_u$
is exponentially small.  Denominator gamma factors give zeros.
\end{enumerate}
There are no other poles between the two lines.  By
\Cref{lem:P-left-line}, the new line for fixed $h,k$ is bounded by
\begin{equation}\label{eq:permutation-new-line}
 T(\log T)^{C_M}(hk)^\eps
 \left\{(T/Y)^\sigma+(h/Y)^\sigma\right\}.
\end{equation}
After the outer sum, choose $M$ so that
\begin{equation}\label{eq:M-choice}
 (B_0-1)\sigma>\frac12(1-\eta)+D+1+2\eps.
\end{equation}
Then \eqref{eq:permutation-new-line} is $O(T^{-D})$ after the outer sum.

\subsection{Joint removability at shift collisions}

The next lemma and proposition are grouped here because they describe the
collision geometry used in the final contour moves.  Logically, however,
they precede \Cref{prop:full-kernel-continuation}.  The lemma is formal local
algebra.  The proposition uses only \eqref{eq:P-lambda}, Euler-factor
symmetry, and $X_{-\gamma,\gamma}=1$; the separate fixed-$s$ application also
uses the residue pairings \eqref{eq:first-common-Z}--\eqref{eq:third-common-Z}.
None of these inputs imports quantitative continuation from
\Cref{prop:full-kernel-continuation}.  This is the forward dependency
recorded explicitly just before that proposition.

\begin{lemma}[Three-divisor cancellation]\label{lem:three-divisor}
Let $x,y$ be independent local holomorphic coordinates and put $z=x+y$.
Suppose $a,b,c,H_0$ are holomorphic and
\begin{equation}\label{eq:three-divisor-form}
 F=\frac{a}{xz}+\frac{b}{yz}+\frac{c}{xy}+H_0,\qquad
 a+c\in(x),\quad b+c\in(y),\quad a-b\in(z).
\end{equation}
Then $F$ is holomorphic across $xyz=0$.
\end{lemma}

\begin{proof}
Over the common denominator, the numerator is
\begin{equation}\label{eq:three-divisor-numerator}
                              N=ay+bx+cz.
\end{equation}
Modulo $x$, use $z=y$ to obtain $N=(a+c)y$; modulo $y$, use $z=x$ to
obtain $N=(b+c)x$; modulo $z$, use $y=-x$ to obtain $N=(b-a)x$.
Thus each of the pairwise coprime linear factors $x,y,z$ divides $N$, and
therefore $xyz\mid N$ in the local analytic ring.
\end{proof}

For the fixed-$s$ assembly, take the coordinates in
\eqref{eq:central-divisors}.  After extracting the simple polar zeta factors,
the direct diagonal, reflected dual diagonal, and zero mode have denominator
patterns $xz$, $yz$, and $xy$.  The three residue cancellations in
\eqref{eq:first-common-Z}--\eqref{eq:third-common-Z} give exactly the three
divisibilities in \eqref{eq:three-divisor-form}.  Thus the complete fixed-$s$
integrand is jointly holomorphic even at double or triple intersections.  Its
separate $1/s$ factor remains, so the prescribed $s=0$ residue is retained.

For the permutation sum, the collision pairs are
\begin{center}
\begin{tabular}{lll}
\toprule
hyperplane & first Euler product & matching Euler product\\
\midrule
$\lambda=-\gamma$
 &$Z_{-\gamma,\beta,\gamma}$
 &$Z_{-\gamma,\beta,\gamma}$, $X_{-\gamma,\gamma}=1$\\
$\lambda=\beta$
 &$Z_{-\gamma,\beta,-\beta}$
 &$Z_{\beta,-\gamma,-\beta}$, equal by first-shift symmetry\\
$\beta=-\gamma$
 &$Z_{\lambda,-\gamma,\gamma}$
 &$Z_{\lambda,-\gamma,\gamma}$, $X_{-\gamma,\gamma}=1$\\
\bottomrule
\end{tabular}
\end{center}
In each row the polar zeta arguments are opposite local coordinates, so the
residues have opposite signs, including the finite factors at primes dividing
$hk$.

\begin{proposition}\label{prop:P-holomorphic}
The complete expression $\Pcal_\lambda(h,k;t)$ in \eqref{eq:P-lambda} is
jointly holomorphic across all shift-collision divisors.
\end{proposition}

\begin{proof}
Apply \Cref{lem:three-divisor} with
\begin{equation}\label{eq:P-collision-coordinates}
 x=\lambda+\gamma,\qquad
 y=-\beta-\gamma,\qquad
 z=\lambda-\beta=x+y.
\end{equation}
The middle, last, and first terms of \eqref{eq:P-lambda} supply the
denominator patterns $xz$, $yz$, and $xy$, respectively.  The three rows of
the preceding table are precisely the required divisibilities.  This also
handles the common collision locus $\lambda=\beta=-\gamma$, rather than only
generic points of the individual hyperplanes.
\end{proof}

\subsection{Proof of the main theorem}

By \eqref{eq:zeta-star-recovery}, the coupled smoothed expansion differs from
\eqref{eq:I-def} by an arbitrary negative power.  Its $r=0$ diagonals and
fixed-orientation zero modes satisfy
\eqref{eq:full-continuation-result}.  The diagonal boundary is reduced by
\eqref{eq:boundary-move-result} and bounded by
\eqref{eq:boundary-zero-bound}.  The complete permutation line contributes
only its $u=0$ residue, by \eqref{eq:permutation-new-line}.  Finally, the
nonzero frequencies are bounded by \eqref{eq:nonzero-final}.  We obtain
\begin{align*}
 \I_{\alpha,\beta,\gamma}(a,b)
 ={}&\sum_{h\leq H}\sum_{k\leq K}
 \frac{a_h\overline{b_k}}{\sqrt{hk}}
 \int_{\R}w(t)\Pcal_\alpha(h,k;t)\dd t\\
 &+O_{\eta,C,\eps}\!\left(
 T^\eps\{KH^{3/2}+T^{1/2}(HK)^{1/2}\}\right).
\end{align*}
By \Cref{prop:P-holomorphic}, the main term is jointly holomorphic at every
shift collision.  This is \eqref{eq:main-asymptotic}.  The replacement
\eqref{eq:X-stirling} follows uniformly from Stirling's formula, completing
the proof.
